\subsection{Условия выпуклости функций на промежутке. Точки перегиба}

\subsubsection*{1. Определения выпуклости}

Лектор отмечает, что в математическом анализе под термином «выпуклая функция» обычно понимают функцию, выпуклую вниз. Выпуклость вверх рассматривается аналогично через умножение функции на $-1$.

\begin{definition}
Функция $f(x)$ называется \textbf{выпуклой вниз} на промежутке $(a, b)$, если для любых двух точек $x_1, x_2 \in (a, b)$ график функции на интервале $(x_1, x_2)$ лежит \textbf{ниже} (в случае строгой выпуклости) или \textbf{не выше} (в случае нестрогой выпуклости) прямой (секущей), проходящей через точки $(x_1, f(x_1))$ и $(x_2, f(x_2))$.
\end{definition}

Аналитически это условие можно записать через соотношение угловых коэффициентов секущих (определение $1'$). Для любых $x_1 < x < x_2$ из интервала $(a, b)$ должно выполняться неравенство:
\[ \frac{f(x) - f(x_1)}{x - x_1} \le \frac{f(x_2) - f(x)}{x_2 - x} \]
Геометрически это означает, что угловой коэффициент секущей, проведенной через точки $x_1$ и $x$, не превосходит углового коэффициента секущей, проведенной через точки $x$ и $x_2$.

\subsubsection*{2. Критерии выпуклости через производные}

\begin{theorem}[Критерий выпуклости по первой производной]
Пусть функция $f(x)$ дифференцируема на интервале $(a, b)$. Функция $f(x)$ является выпуклой вниз тогда и только тогда, когда её производная $f'(x)$ не убывает на этом интервале. Для строгой выпуклости необходимо и достаточно строгого возрастания $f'(x)$.
\end{theorem}

\begin{tcolorbox}[title=Доказательство, breakable]
\textbf{1. Необходимость ($\implies$):}
Пусть функция выпукла вниз. Используем аналитическое условие для $x_1 < x < x_2$:
\[ \frac{f(x) - f(x_1)}{x - x_1} \le \frac{f(x_2) - f(x)}{x_2 - x} \]
Устремим $x \to x_1 + 0$ в левой части. По определению производной, предел разностного отношения равен $f'(x_1)$. Получаем:
\[ f'(x_1) \le \frac{f(x_2) - f(x_1)}{x_2 - x_1} \]
Теперь в исходном неравенстве устремим $x \to x_2 - 0$ в правой части. Предел правой части равен $f'(x_2)$. Получаем:
\[ \frac{f(x_2) - f(x_1)}{x_2 - x_1} \le f'(x_2) \]
Объединяя эти неравенства, имеем $f'(x_1) \le f'(x_2)$ для любых $x_1 < x_2$, что означает неубывание производной.

\textbf{2. Достаточность ($\impliedby$):}
Пусть $f'(x)$ не убывает. Рассмотрим произвольные точки $x_1 < x < x_2$. Применим теорему Лагранжа к интервалам $[x_1, x]$ и $[x, x_2]$:
\[ \frac{f(x) - f(x_1)}{x - x_1} = f'(\xi_1), \quad \xi_1 \in (x_1, x) \]
\[ \frac{f(x_2) - f(x)}{x_2 - x} = f'(\xi_2), \quad \xi_2 \in (x, x_2) \]
Так как $x_1 < \xi_1 < x < \xi_2 < x_2$, то $\xi_1 < \xi_2$. В силу неубывания производной $f'(\xi_1) \le f'(\xi_2)$, откуда следует выполнение аналитического условия выпуклости.
\end{tcolorbox}

\begin{theorem}[Критерий выпуклости по второй производной]
Пусть функция $f(x)$ дважды дифференцируема на интервале $(a, b)$. Тогда:
\begin{enumerate}
    \item $f(x)$ выпукла вниз на $(a, b) \iff f''(x) \ge 0$ на $(a, b)$.
    \item Если $f''(x) > 0$ на $(a, b)$, то функция $f(x)$ строго выпукла вниз на этом интервале.
\end{enumerate}
\end{theorem}

\begin{tcolorbox}[title=Доказательство, breakable]
Утверждение является прямым следствием предыдущей теоремы и условий монотонности функции. Производная $f'(x)$ не убывает тогда и только тогда, когда её производная (то есть $f''(x)$) неотрицательна. Соответственно, если $f''(x) > 0$, то $f'(x)$ строго возрастает, что гарантирует строгую выпуклость.
\end{tcolorbox}

\subsubsection*{3. Точки перегиба}

Лектор подчеркивает, что определение точки перегиба только как точки смены направления выпуклости является недостаточным.

\begin{definition}
Точка $x_0$ называется \textbf{точкой перегиба} графика функции $f(x)$, если:
\begin{enumerate}
    \item В этой точке существует \textbf{касательная} к графику функции (возможно, вертикальная).
    \item При переходе через точку $x_0$ направление выпуклости функции меняется на противоположное.
\end{enumerate}
\end{definition}

\textbf{Геометрический смысл:} В точке перегиба график функции переходит с одной стороны касательной на другую. Если функция выпукла вниз, она лежит выше своей касательной; если выпукла вверх — ниже. В точке перегиба происходит смена этого расположения.

\begin{theorem}[Необходимое условие точки перегиба]
Если $x_0$ является точкой перегиба функции $f(x)$ и в этой точке существует вторая производная $f''(x_0)$, то $f''(x_0) = 0$.
\end{theorem}

\begin{theorem}[Достаточное условие точки перегиба]
Пусть функция $f(x)$ непрерывна в точке $x_0$ и дважды дифференцируема в её проколотой окрестности. Если при переходе через $x_0$ вторая производная $f''(x)$ меняет знак, то $x_0$ является точкой перегиба.
\end{theorem}

%\includegraphics[width=0.5\textwidth]{inflection_point.png}
% На рисунке должен быть изображен график функции (например, y=x^3), 
% касательная в точке (0,0) и отмечено, что до нуля функция выпукла вверх, 
% а после — выпукла вниз, при этом график пересекает касательную.